\documentclass[12pt,a4paper]{article}

\usepackage{amsmath, amssymb, amsthm}
\usepackage{geometry}
\usepackage{hyperref}
\usepackage{algorithm}
\usepackage{algpseudocode}
\usepackage{booktabs}
\usepackage{xcolor}
\usepackage{enumitem}
\usepackage{microtype}

\geometry{margin=2.5cm}

\hypersetup{
  colorlinks=true,
  linkcolor=blue!60!black,
  urlcolor=blue!60!black,
}

\newtheorem{definition}{Definition}
\newtheorem{proposition}{Proposition}
\newtheorem{remark}{Remark}

\title{\textbf{Ghost Font: A Formal Specification of\\
Motion-Encoded Anti-AI Typography}}
\author{
  \textit{Formalized from the work of Eric Lu / Mixfont (2026)}\\[0.4em]
  \small Research analysis and mathematical specification
}
\date{July 2026}

\begin{document}

\maketitle
\tableofcontents
\newpage

%──────────────────────────────────────────────────────────────────────────────
\section{Introduction}
%──────────────────────────────────────────────────────────────────────────────

Ghost Font, released by Eric Lu and the Mixfont team in July 2026, is a
motion-based visual encoding system that embeds a short text message into a
video clip in a form immediately legible to human observers but opaque to
current AI multimodal systems.  Unlike traditional steganographic schemes
that hide information in the spatial or spectral domain of a single frame,
Ghost Font encodes information in the \emph{temporal domain}—specifically in
the relative motion direction of two overlapping dot populations.  This
document provides a formal mathematical specification of that encoding
algorithm.

\subsection{Motivation and Context}

Contemporary large language models with vision capabilities process video by
decomposing it into a sequence of individual frames and applying image-based
analysis to each frame independently.  They do not, in general, perform
native temporal integration of motion across frames.  The human visual
system, by contrast, possesses dedicated neural machinery (the dorsal
``where/how'' stream, area MT/V5) for detecting and integrating motion across
time.  Ghost Font exploits this architectural asymmetry: the letter signal
exists only in the \emph{inter-frame motion differential}, which humans
perceive instantly but which current AI frame-by-frame pipelines cannot
recover without explicit motion-compensation preprocessing.

A secondary layer—a \emph{decoy message}—is superimposed as a faint static
pattern.  When an AI system averages frames or scans for spatial patterns, it
surfaces this decoy rather than the true message, acting as a honeypot.

%──────────────────────────────────────────────────────────────────────────────
\section{Notation and Primitives}
%──────────────────────────────────────────────────────────────────────────────

\begin{definition}[Canvas]
A \emph{canvas} $\mathcal{C} = (W, H)$ is a discrete integer grid of pixels
with width $W \in \mathbb{Z}^+$ and height $H \in \mathbb{Z}^+$.  A pixel is
addressed by $(x, y)$ with $x \in [0, W)$, $y \in [0, H)$, all integers.
\end{definition}

\begin{definition}[Frame]
A \emph{frame} at time step $t$ is a function
$F_t : \mathcal{C} \to \{0, 1\}$, where $0$ denotes a background (light)
pixel and $1$ denotes a foreground (dark) dot.
\end{definition}

\begin{definition}[Dot Population]
A \emph{dot population} $\mathcal{P}$ is a finite multiset of particles.
Each particle $p_i$ carries an initial position
$(x_i^{(0)}, y_i^{(0)}) \in \mathcal{C}$ drawn uniformly at random, and
belongs to exactly one of two classes: \emph{signal} or \emph{background}.
\end{definition}

\begin{definition}[Letter Mask]
The \emph{letter mask} $\mathcal{M} : \mathcal{C} \to \{0, 1\}$ is a binary
function derived by rasterizing the message glyphs onto the canvas at a
chosen font face and size.  $\mathcal{M}(x, y) = 1$ iff pixel $(x, y)$ falls
within a glyph region.
\end{definition}

%──────────────────────────────────────────────────────────────────────────────
\section{The Ghost Font Encoding Algorithm}
%──────────────────────────────────────────────────────────────────────────────

The algorithm proceeds in four stages: (1)~canvas and mask initialization,
(2)~dual-population dot seeding, (3)~frame generation via opposing motion,
and (4)~decoy superimposition.

\subsection{Stage 1: Canvas and Mask Initialization}

Given a message string $s$, rasterize each glyph onto the canvas using a
chosen typeface, producing the letter mask $\mathcal{M}$.  Glyphs are
typically laid out horizontally centred on the canvas.

\subsection{Stage 2: Dual-Population Dot Seeding}

Let $\rho \in (0, 1)$ be the target dot density.  Seed two disjoint
populations uniformly at random:
\begin{align}
  \mathcal{P}_{\mathrm{sig}} &= \bigl\{ p_i :
    (x_i^{(0)}, y_i^{(0)}) \sim \mathrm{Uniform}(\mathcal{C}),\;
    \mathcal{M}(x_i^{(0)}, y_i^{(0)}) = 1 \bigr\}, \\
  \mathcal{P}_{\mathrm{bg}}  &= \bigl\{ p_j :
    (x_j^{(0)}, y_j^{(0)}) \sim \mathrm{Uniform}(\mathcal{C}),\;
    \mathcal{M}(x_j^{(0)}, y_j^{(0)}) = 0 \bigr\},
\end{align}
with $|\mathcal{P}_{\mathrm{sig}}| + |\mathcal{P}_{\mathrm{bg}}|
\approx \rho \cdot W \cdot H$.

\begin{remark}
The mask determines \emph{class membership}, not ongoing spatial
confinement.  At time $t > 0$, particles move freely across the entire
canvas; a signal particle originally inside a glyph region may later be
anywhere.
\end{remark}

\subsection{Stage 3: Frame Generation via Opposing Motion}

The encoding mechanism assigns opposite vertical velocity vectors to the
two populations.

\begin{definition}[Signal and Background Velocities]
Let $v \in \mathbb{R}^+$ be the nominal scroll speed in pixels per second,
$\Delta t = 1/f_r$ the inter-frame interval for frame rate $f_r$, and
$(w_x, w_y) \in \mathbb{R}^2_{\geq 0}$ optional wobble amplitudes with
angular frequencies $(\omega_x, \omega_y)$.  Then
\begin{align}
  \mathbf{v}_{\mathrm{sig}}(n) &=
    \bigl(w_x \sin(\omega_x n\,\Delta t),\;
           -v + w_y \cos(\omega_y n\,\Delta t)\bigr)
    \quad (\text{upward}), \\
  \mathbf{v}_{\mathrm{bg}} &= (0,\; +v)
    \quad (\text{downward, constant}).
\end{align}
(Upward is the $-y$ direction in standard image coordinates.)
\end{definition}

The position of signal particle $p_i$ at integer frame index $n$ is
\begin{equation}
  \begin{pmatrix} x_i(n) \\ y_i(n) \end{pmatrix}
  =
  \begin{pmatrix}
    x_i^{(0)} + \displaystyle\sum_{k=0}^{n-1} w_x \sin(\omega_x k\,\Delta t) \\[6pt]
    \Bigl(y_i^{(0)} - n\,v\,\Delta t
          + \displaystyle\sum_{k=0}^{n-1} w_y \cos(\omega_y k\,\Delta t)\Bigr)
    \bmod H
  \end{pmatrix},
\end{equation}
with vertical wrap-around modulo $H$ to maintain uniform coverage.
The background particle $p_j$ satisfies
\begin{equation}
  (x_j(n), y_j(n)) = (x_j^{(0)},\; (y_j^{(0)} + n\,v\,\Delta t) \bmod H).
\end{equation}

Each frame $F_n$ is rendered by placing a filled disk of radius $r$ at each
particle position:
\begin{equation}
  F_n(x, y) = 1 \iff
  \exists\, p_i \in \mathcal{P}_{\mathrm{sig}} \cup \mathcal{P}_{\mathrm{bg}}
  \text{ s.t. }
  \bigl\|(x, y) - (x_i(n), y_i(n))\bigr\|_2 \leq r.
\end{equation}

\subsection{Stage 4: Decoy Superimposition}

Let $\mathcal{D} : \mathcal{C} \to \{0,1\}$ be the binary raster of a decoy
string $d$ (e.g., \texttt{WRITTEN IN GHOST FONT}).  Each output frame is
\begin{equation}
  F_n^{\mathrm{out}}(x, y)
  = \min\!\bigl(1,\; F_n(x, y) + \alpha_d\,\mathcal{D}(x, y)\bigr),
\end{equation}
where $\alpha_d \ll 1$ is the decoy opacity, chosen so that the decoy is
invisible in any single frame but accumulates to visibility when frames are
temporally averaged (see Section~\ref{sec:security}).

\begin{algorithm}[H]
\caption{Ghost Font Video Encoding}
\begin{algorithmic}[1]
\Require Message $s$, canvas $(W,H)$, speed $v$, frame rate $f_r$,
  dot density $\rho$, wobble $(w_x,w_y,\omega_x,\omega_y)$,
  total frames $N_f$, decoy string $d$, decoy opacity $\alpha_d$
\Ensure Video $\mathbf{V} = \bigl[F_0^{\mathrm{out}},\ldots,F_{N_f-1}^{\mathrm{out}}\bigr]$
\State $\mathcal{M} \leftarrow \textsc{RasterizeGlyphs}(s,\,W,\,H)$
\State $\mathcal{D} \leftarrow \textsc{RasterizeGlyphs}(d,\,W,\,H)$
\State $\mathcal{P}_{\mathrm{sig}},\,\mathcal{P}_{\mathrm{bg}} \leftarrow \textsc{SeedDots}(\mathcal{M},\,\rho,\,W,\,H)$
\For{$n = 0$ \textbf{to} $N_f - 1$}
  \State $F_n \leftarrow \mathbf{0}_{W \times H}$
  \For{each $p_i \in \mathcal{P}_{\mathrm{sig}}$}
    \State $(x,y) \leftarrow \textsc{AdvanceSig}(p_i,\,n,\,v,\,\Delta t,\,w_x,\,w_y,\,\omega_x,\,\omega_y,\,H)$
    \State $\textsc{PlaceDot}(F_n,\,x,\,y,\,r)$
  \EndFor
  \For{each $p_j \in \mathcal{P}_{\mathrm{bg}}$}
    \State $(x,y) \leftarrow \textsc{AdvanceBg}(p_j,\,n,\,v,\,\Delta t,\,H)$
    \State $\textsc{PlaceDot}(F_n,\,x,\,y,\,r)$
  \EndFor
  \State $F_n^{\mathrm{out}} \leftarrow \min(1,\;F_n + \alpha_d\cdot\mathcal{D})$
\EndFor
\Return $\mathbf{V}$
\end{algorithmic}
\end{algorithm}

%──────────────────────────────────────────────────────────────────────────────
\section{Perceptual Properties}
%──────────────────────────────────────────────────────────────────────────────

\subsection{Human Legibility via Kinetic Depth and Motion Segregation}

Human motion perception (\emph{kinetic depth effect}, \emph{structure-from-motion})
allows the visual system to segment two coherently moving dot fields even when
they overlap and each individually appears random.  The signal region—dots
scrolling upward—segregates from the background—dots scrolling downward—
making letter shapes immediately apparent to a human viewer, typically within
one second of video playback.

\subsection{Single-Frame Uninformativeness}

\begin{proposition}
\label{prop:uninformative}
For any single frame $F_n$ produced by Algorithm~1 (ignoring the
$\alpha_d$ decoy term), the pixel distribution is i.i.d.\
$\mathrm{Bernoulli}(\rho)$, independent of the letter mask $\mathcal{M}$.
\end{proposition}

\begin{proof}[Proof sketch]
Each particle's initial position is drawn uniformly from $\mathcal{C}$.
Advancing by a deterministic offset (modulo $H$) maps a uniform distribution
to a uniform distribution.  The union rendering of two independent
uniform-density Bernoulli processes at combined rate $\rho$ is again
Bernoulli($\rho$) at each pixel, with no spatial structure that depends on
which population placed the dot.
\end{proof}

Consequently, a screenshot (single frame) carries zero mutual information
about the letter mask.

\subsection{Message Recovery by Motion-Compensated Integration}
\label{sec:recovery}

Define the vertical shift operator $\mathcal{T}_{\delta y}[F]$ that
translates frame $F$ upward by $\delta y$ pixels (with wrap-around).
The signal population at frame $n$ is displaced by
$\delta y_n = n\,v\,\Delta t$ (ignoring wobble) relative to frame $0$.
Applying the inverse shift before summation yields
\begin{equation}
  I_{\mathrm{sig}}(x, y)
  = \frac{1}{N_f} \sum_{n=0}^{N_f-1}
      \mathcal{T}_{n\,v\,\Delta t}\!\bigl[F_n\bigr](x, y).
\end{equation}
In this shifted coordinate, signal dots accumulate coherently while
background dots are smeared uniformly across all $y$-positions.  The
result is a high-contrast image of the letter mask $\mathcal{M}$, with
signal-to-noise ratio growing as $\mathcal{O}(\sqrt{N_f})$.

%──────────────────────────────────────────────────────────────────────────────
\section{Security Analysis}
\label{sec:security}
%──────────────────────────────────────────────────────────────────────────────

\subsection{Layer 1 — Single-Frame Opacity}

By Proposition~\ref{prop:uninformative}, no information about $s$ is
recoverable from any individual frame.  Image-based AI analysis (screenshot,
single-frame OCR) yields null information about the true message.

\subsection{Layer 2 — Temporal Averaging and the Decoy Honeypot}

Na\"{\i}ve temporal averaging without motion compensation gives
\begin{equation}
  \bar{F}(x,y)
  = \frac{1}{N_f}\sum_{n=0}^{N_f-1} F_n^{\mathrm{out}}(x,y)
  \approx \rho \cdot \mathbf{1} + \alpha_d\,\mathcal{D}(x,y).
\end{equation}
The moving populations average out to uniform density $\rho$; the only
spatially structured component that survives is the \emph{static} decoy
$\alpha_d\,\mathcal{D}$.  An AI applying frame averaging therefore recovers
$d$, not $s$—the honeypot property.  The model reports this finding with
high confidence, further masking the existence of the true message.

\subsection{Layer 3 — Motion-Compensation Break}

A determined adversary with local code execution can defeat both layers:
\begin{enumerate}[label=\arabic*.]
  \item Estimate the scroll velocity $v$ via optical flow (e.g.,
    Farneb\"{a}ck or Lucas--Kanade).
  \item Apply the motion-compensated integrator of
    Section~\ref{sec:recovery} to reconstruct $\mathcal{M}$.
  \item Run OCR on the reconstructed image.
\end{enumerate}
This constitutes a complete break.  The \texttt{ghost-font-breaker} tool
(42-evey, GitHub, 2026) demonstrates deterministic message recovery on all
publicly available Ghost Font clips with recovered velocity
$v = 4.012\,\mathrm{px/frame}$, wobble amplitudes
$w_x = 72\,\mathrm{px}$, $w_y = 48\,\mathrm{px}$.

\subsection{Resistance Summary}

\begin{center}
\renewcommand{\arraystretch}{1.3}
\begin{tabular}{p{4.5cm}p{4cm}cc}
\toprule
\textbf{Adversary capability} & \textbf{Method} &
  \textbf{Decoy?} & \textbf{True msg?} \\
\midrule
Screenshot / single frame    & Direct OCR           & No  & No  \\
Frame averaging              & Na\"{\i}ve avg.      & Yes & No  \\
Online multimodal model      & Frame-by-frame       & Yes & No  \\
Local code execution         & Motion-compensated avg. & Yes & Yes \\
\bottomrule
\end{tabular}
\end{center}

%──────────────────────────────────────────────────────────────────────────────
\section{Parameter Reference}
%──────────────────────────────────────────────────────────────────────────────

\begin{table}[h!]
\centering
\caption{Ghost Font algorithm parameters (observed / recovered values)}
\label{tab:params}
\renewcommand{\arraystretch}{1.2}
\begin{tabular}{llll}
\toprule
\textbf{Parameter} & \textbf{Symbol} & \textbf{Observed value} & \textbf{Notes} \\
\midrule
Scroll speed           & $v$                          & 120\,px/s (4.012\,px/frame) & at $\approx$30\,fps \\
Signal direction       & $\mathbf{v}_{\mathrm{sig}}$  & upward ($-y$)               & \\
Background direction   & $\mathbf{v}_{\mathrm{bg}}$   & downward ($+y$)             & \\
Wobble amplitude $x$   & $w_x$                        & 72\,px                      & sinusoidal \\
Wobble amplitude $y$   & $w_y$                        & 48\,px                      & cosinusoidal \\
Dot density            & $\rho$                       & $\approx 0.3$--$0.5$        & \\
Dot radius             & $r$                          & 2--3\,px                    & \\
Decoy opacity          & $\alpha_d$                   & $\ll 1$                     & survives avg. \\
Decoy string           & $d$                          & \texttt{WRITTEN IN GHOST FONT} & hardcoded \\
Max message length     & $|s|$                        & 36 characters               & prototype \\
\bottomrule
\end{tabular}
\end{table}

%──────────────────────────────────────────────────────────────────────────────
\section{Limitations and Future Directions}
%──────────────────────────────────────────────────────────────────────────────

\begin{enumerate}[label=\arabic*.]
  \item \textbf{Human legibility cost.}  Ghost Font is harder for humans to
    read than ordinary text.  The trade-off is inherent: motion encoding
    requires a video medium and reduces reading speed compared to static
    typography.

  \item \textbf{Video-native AI models.}  The current resistance relies on AI
    systems processing video as frame sequences.  Once video-native models
    performing native temporal integration of motion become commonplace, the
    core evasion mechanism will likely fail.

  \item \textbf{No cryptographic security.}  Ghost Font provides obfuscation,
    not cryptographic secrecy.  A motivated adversary with code execution can
    fully recover the message.  True secrecy requires a shared key; a
    companion design document proposes key-encrypted extensions.

  \item \textbf{Character limit.}  The current prototype supports short
    messages ($\leq 36$ characters).  Longer messages reduce per-letter
    motion contrast and degrade human legibility.

  \item \textbf{CAPTCHA potential.}  Motion-based encoding is a promising
    direction for next-generation CAPTCHAs, exploiting the human advantage in
    biological motion perception where static-image CAPTCHAs have been largely
    defeated by current vision models.
\end{enumerate}

%──────────────────────────────────────────────────────────────────────────────
\section{Conclusion}
%──────────────────────────────────────────────────────────────────────────────

Ghost Font formalizes as a \emph{dual-population, opposing-direction random
dot kinematogram} in which the letter mask $\mathcal{M}$ partitions dots into
a signal population (scrolling up at velocity $v$) and a background population
(scrolling down at $v$).  Single frames are statistically uniform
(Proposition~\ref{prop:uninformative}); the letter mask is recoverable only
by motion-compensated temporal integration (Section~\ref{sec:recovery}).  A
static decoy layer provides a honeypot for frame-averaging adversaries.  The
system is provably opaque to single-frame analysis, practically opaque to
current online multimodal AI, and fully defeatable by local code-execution via
optical flow and motion-compensated integration—a known limitation that
motivates the key-encrypted extensions described in the companion design
document.

\end{document}
